\documentclass[11pt]{article}
\usepackage[margin=0.9in]{geometry}
\usepackage{amsmath,amssymb,amsthm,booktabs,mathrsfs}
\usepackage[T1]{fontenc}
\usepackage{lmodern,microtype}
\usepackage[colorlinks=true,linkcolor=blue,urlcolor=blue]{hyperref}
\newtheorem{theorem}{Theorem}[section]
\newtheorem{lemma}[theorem]{Lemma}
\newtheorem{proposition}[theorem]{Proposition}
\newtheorem{corollary}[theorem]{Corollary}
\theoremstyle{remark}
\newtheorem{remark}[theorem]{Remark}
\newcommand{\norm}[1]{\left\|#1\right\|}
\title{Toward a four-dimensional Yang--Mills construction\\
\large Notebook 31: finite local induction and coupling extraction}
\author{Research working document}
\date{17 October 2026}
\begin{document}
\maketitle
\begin{abstract}
We carry the renormalized homogeneous local induction to actual
finite periodic regulators. A short-lift Maxwell extraction is
normalized by the actual finite classical action. The finite Ward
defects are repaired degree by degree; periodic summation by parts
and connected localization give a classical contraction with an
explicit exponentially small wrapping error. The carrier exponent
and field radius are unchanged.

The finite quantum map descends directly to equality classes of
normalized total interactions. Its input Lipschitz bound, its
quadratic zero-input force, and the finite classical contraction
give a uniform block induction. At every intermediate stage an
exactly equivalent connected family lies within the next integral's
strict budget. The change of representative introduces no
approximation error; all raw constants are compensated in the
vacuum factor. The number of permitted finite blocks does not
enter the amplitude constants.

We also prove evenness in the coupling at fixed retained threshold.
A uniform complex coupling disk and Cauchy estimates identify the
leading zero-input correction as the pivot term plus the Gaussian
determinant, with error bounded by
\(C[\tau(T)+g^2T^6]\). The physical logarithmic threshold is substituted
only after this two-parameter estimate is established.

On the jointly compatible infinite interaction space, a
quadratic-only spatial estimate at fixed field radius and the
classical semigroup prove exact preservation of the Maxwell
coefficient. Thus the coupling drift on an extracted input
comes entirely from the quantum remainder. This infinite
identity does not set finite cutoff-extraction errors to zero.

These results concern the specified local macrostep and its common
analytic continuation. The physical inverse coupling must satisfy
the exact extraction recurrence and remain in the weak range.
Its drift sign and ultraviolet tuning, the other field sectors,
and the continuum observable construction remain to be proved.
\end{abstract}

The goal remains a nontrivial four-dimensional Yang--Mills quantum
theory. This notebook uses the actual finite classical and quantum
constructions of Notebooks 15--30, with their fixed gauge, retained
window, trace convention, and common domains. It closes the finite
representative-budget issue left explicit in Notebook 30.

